\documentclass{ejercicios}

\begin{document}

\tableofcontents
\clearpage

\begin{exercise}
  Sea $\varphi\in \mathcal{C}([a,b],\mathbb{R})$ tal que $\varphi(a) = x_a$ y $\varphi(b) = x_b$. Determine una función $\hat{x}\in\hat{\mathcal{D}} = \{\hat{x}\in\mathcal{P}\mathcal{C}^1([a,b],\mathbb{R}):\hat{x}(a) = x_a$ y $\hat{x}(b) = x_b\}$ con la propiedad de que la región $\Omega$ limitada por las curvas de las gráficas de las funciones $\varphi$ y $\hat{x}$ sea máxima, tomando en consideración el supuesto de que la longitud de la curva $\hat{x}$ es una constante $l>0$.

\end{exercise}

\begin{solution}
  Primero identifiquemos que el funcional que vamos a trabajar esla dado en realidad por 
  \[\Omega[\hat{x}] = \int_{a}^{b}\rvert\varphi(t) - \hat{x}(t)\lvert\,dt\]
  además como el problema nos dice que la longitud de la curva $\hat{x}$ es una constante $l>0$ entonces tenemos la condición 
  \[G[x] = \int_{a}^{b}\left(\sqrt{1 + (\dot{x}')^2} - \frac{l}{b-a}\right)\,dt = 0\]
  asi con esto en mente podemos buscar los valores máximos de 
  \[
  \begin{aligned}
    & \hat{F}(\hat{x}) = \lambda_0\Omega[{\hat{x}}] + \lambda_1 G[\hat{x}]\\ 
    & = \int_{a}^{b}\left(\lambda_0\lvert\varphi - \hat{x}\rvert + \lambda_1\sqrt{1 + (\hat{x}')^2} - \frac{l}{b-a}\right)\,dt
  \end{aligned}
  \]
  e identifiquemos $\hat{L} = \lambda_0\lvert\varphi - \hat{x}\rvert + \lambda_1\sqrt{1 + (\hat{x}')^2} - \frac{l}{b-a}$ ahora de la nueva \textbf{ecuación de Euler-Lagrange} se tiene que 
  \[
    0 = \frac{\partial\hat{L}}{\partial{\hat{x}}}-\frac{d}{\,dt}\frac{\partial\hat{L}}{\partial{\dot{\hat{x}}}} =  \lambda_0sign(\varphi - \hat{x}) - \lambda_1\frac{d}{\,dt}\left(\frac{\dot{\hat{x}}}{\sqrt{1 + \dot{\hat{x}}}}\right)
  \]
  de aqui si $\lambda_1 \neq 0$ definamos $\lambda = \frac{\lambda_0}{\lambda_1}$
\[\frac{\dot{\hat{x}}}{\sqrt{1 + \dot{\hat{x}}}} = \lambda\int sign(\varphi - \hat{x})\,dt\]
  ahora observemos que 
  \[
    \int sign(\varphi - \hat{x})\,dt = 
  \begin{cases}
    1 & \text{si } \varphi - \hat{x} > 0 \\ 
    0 & \text{si } \varphi - \hat{x} = 0 \\ 
    -1 & \text{si } \varphi - \hat{x} < 0
  \end{cases}
  \]
  veamos los casos 
  \begin{enumerate}
    \item $\varphi > \hat{x}$ \\ Entonces \[\frac{\dot{\hat{x}}}{\sqrt{1 + \dot{\hat{x}}}} = \lambda\] de esto se sigue 
      \[\dot{\hat{x}}^2 = \lambda^2(1 + \dot{\hat{x}})\]
      \[\dot{\hat{x}}^2 - \lambda^2\dot{\hat{x}} - \lambda^2 = 0\]

      de aqui podemos ver que la ecuación tiene forma polinomial $r^2 - C_1 r - C_1 = 0$ el cual cuenta con soluciones 
    \[ r_{1,2} = \frac{C_1 \pm \sqrt{{C_1}^2 + 4C_1}}{2}\]
  \end{enumerate}
\end{solution}

\begin{exercise}
  Hallar los valores extremos del funcional 
  \[
    J[x] = \int_{0}^{\pi}\dot{x}^2(t)\,dt
  \]
  donde $\mathcal{D} = \{x\in\mathcal{C}^1([a,b],\mathbb{R}):x(0) = x(\pi) = 0$ y $G[x] = 1\}$ y 
  \[
    G[x] = \int_{0}^{\pi}\dot{x}^2(t)\,dt
  \]
\end{exercise}

\begin{solution}
  Primero consideremos 
\[\hat{G}[x] = \int_{0}^{\pi} (\dot{x}^2(t) - 1)\,dt\]
con esto podemos considerar 
\[\hat{F}(x) = \int_{0}^{\pi} ((\lambda_0 - \lambda_1)\dot{x}^2(t) - 1)\,dt \]
y sea $\hat{L}(x) = (\lambda_0 - \lambda_1)\dot{x}^2(t) - 1$ entonces la ecuación de \textbf{Euler-Lagrange} esta dada por
\[0 = \frac{\partial}{\partial x}\hat{L} - \frac{d}{\,dt}\frac{\partial}{\partial\dot{x}}\hat{L}\]
identifiquemos que \[\frac{\partial}{\partial x}\hat{L} = 0\] y 
\[\frac{d}{\,dt}\frac{\partial}{\partial\dot{x}}\hat{L} = \frac{d}{\,dt}(2(\lambda_0 - \lambda_1)\dot{x})\]
entonces
\[\frac{d}{\,dt}(2(\lambda_0 - \lambda_1)\dot{x}) = 0\]
de donde se sigue que 
\[2(\lambda_0 - \lambda_1)\dot{x} = k\]
con $k\in\mathbb{R}$ por lo tanto $x(t) = \frac{k}{2(\lambda_0 - \lambda_1)}t + C_1$
\end{solution}

\begin{exercise}
  Sean $p$, $q$, $r:[a,b]\to\mathbb{R}$ funciones continuas y sea $J:\mathcal{D}\to\mathcal{R}$ el funcional
  \[J[x] = \int_{a}^{b}(p(t)\dot{x}^2(t) - q(t)x^2(t))\,dt\]
  donde $\mathcal{D} = \{x\in\mathcal{C}^1([a,b],\mathbb{R}):x(a) = x_a, \dot{x}(b) = x_b$ y $G[x] = 0\}$ con 
  \[G[x] = \int_{a}^{b}r(t)x^2(t)\,dt - 1\]
  Mostrar que los valores extremos del funcional son soluciones de la ecuación diferencial
  \[
    \frac{d}{dt}\big(p(t)\frac{dx}{dt}\big) + (q(t) + \lambda r(t))x = 0
  \]
  donde $\lambda$ es una constante.
\end{exercise}

\begin{solution}
  \noindent
  Como $J$ y $G$ estan dadas, entonces definamos 
  \[\hat{F}[x] = \lambda_0J[x] + \lambda_1G[x] = \int_{a}^{b} \left(\lambda_0\left(p(t)\dot{x}^2(t) - q(t)x^2(t) \right) + \lambda_1\left(r(t)x^2(t) - \frac{1}{b-a}\right)\right)\,dt\]
  de donde identificamos 
  \[\hat{L} = \lambda_0\left(p(t)\dot{x}^2(t) - q(t)x^2(t) \right) + \lambda_1\left(r(t)x^2(t) - \frac{1}{b-a}\right)\]
  con esto la ecuación de \textbf{Euler-Lagrange} esta dada por 
  \begin{equation}\label{eq:1}
  0 = \frac{\partial}{\partial x}\hat{L} - \frac{d}{\,dt}\frac{\partial}{\partial}\dot{x}\hat{L}
  \end{equation}
  con 
  \[
  \begin{aligned}
    & \frac{\partial}{\partial x}\hat{L} = -2\lambda_0q(t)x(t) + 2\lambda_1r(t)x(t)\\ 
    & \frac{\partial}{\partial\dot{x}}\hat{L} = 2\lambda_0p(t)\dot{x}(t)
  \end{aligned}
  \]
  sustituyendo esto último en la ecuación~\ref{eq:1} se tiene que 
  \begin{equation}\label{eq:2}
  0 = -2\left(\lambda_0q(t)x(t) - \lambda_1r(t)x(t) + \frac{d}{\,dt}\left(\lambda_0p(t)\dot{x}(t)\right)\right)
  \end{equation}
  si $\lambda_0\neq 0$ entonces sea $\lambda = \frac{\lambda_1}{\lambda_0}$ entonces la ecuación~\ref{eq:2} se transforma en 
  \[\frac{d}{\,dt}\left(p(t)\frac{dx}{\,dt}\right) + (q(t) - \lambda r(t))x(t) = 0\]
  que es lo que se quería demostrar.
\end{solution}

\begin{exercise}\label{ej:4}
  Sean $L,M\in\mathcal{C}^2([a,b]\times\mathbb{R}^4,\mathbb{R})$ y $J:\mathcal{D}_{x_1}\times\mathcal{D}_{x_2}\to\mathbb{R}$ el funcional definido por
  \[
    J[x_1,x_2] = \int_{a}^{b}L(t,x_1(t),x_2(t),\dot{x_1}(t),\dot{x_2}(t))\,dt,
  \]
  donde
  \[
  \begin{aligned}
    & \mathcal{D}_{x_1}\times\mathcal{D}_{x_2} = \{(x_1,x_2)\in\mathcal{C}^1([a,b],\mathbb{R})\times\mathcal{C}^1([a,b],\mathbb{R}):x_k(a) = {x_a}^k\text{ y }x_k(b) = {x_b}^k \\ 
    & \text{ para } k = 1,2 \text{ y } G[x_1,x_2] = c\},
  \end{aligned}
\]
  con $c$ una constante dada y 
  \[
    G[x_1,x_2] = \int_{a}^{b}M(t,x_1(t),x_2(t),\dot{x_1}(t),\dot{x_2}(t))\,dt .
  \]
  Hallar la condición necesaria de optimalidad.
\end{exercise}

\begin{solution}
  Primero recordemos que un funcional $J[x_1,x_2]$ como el descrito tiene ecuación d Euler-Lagrange de la forma siguiente 
  \begin{equation}\label{eq:3}
  \frac{\partial}{\partial x_k}L - \frac{d}{\,dt}\frac{\partial}{\partial\dot{x_k}}L = 0
  \end{equation}
  en nuestro caso $k\in\{1,2\}$, lo cual permite crear un sistema de ecuaciones; primero vamos a definir 
  \begin{equation}
    \hat{F}[x_1,x_2] = \lambda_0J[x_1,x_2] + \lambda_1\hat{G}[x_1,x_2] = \int_{a}^{b} \left(\lambda_0L(t,x_1,x_2,\dot{x_1},\dot{x_2}) + \lambda_1\left(M(t,x_1,x_2,\dot{x_1},\dot{x_2}) - \frac{c}{b-a}\right)\right)\,dt
  \end{equation}
  e identifiquemos 
  \begin{equation}
    \hat{L} = \lambda_0L(t,x_1,x_2,\dot{x_1},\dot{x_2}) + \lambda_1\left(M(t,x_1,x_2,\dot{x_1},\dot{x_2}) - \frac{c}{b-a}\right)
  \end{equation}
  entonces de la ecuación de Euler-Lagrange descrita antes se sigue que 
  \[
  \begin{aligned}
    & \frac{\partial}{\partial x_k}\hat{L} &= \lambda_0\frac{\partial}{\partial x_k}L + \lambda_1\frac{\partial}{\partial x_k}M \\ 
    & \frac{\partial}{\partial\dot{x_k}}\hat{L} &= \lambda_0\frac{\partial}{\partial\dot{x_k}}L + \lambda_1\frac{\partial}{\partial\dot{x_k}}M
  \end{aligned}
  \]
  ahora si $\lambda_0\neq 0$ definamos $\lambda = \frac{\lambda_1}{\lambda_0}$ y sustituimos lo anterior en la ecuación~\ref{eq:3} se sigue que
  \[
    \frac{\partial}{\partial x_k}L + \lambda\frac{\partial}{\partial x_k}M - \left(\frac{d}{\,dt}\frac{\partial}{\partial\dot{x_k}}L + \lambda\frac{d}{\,dt}\frac{\partial}{\partial\dot{x_k}}M\right) = 0
  \]
  esta seria la condición necesario de optimalidad sujeto a las condiciones 
  \[x_k(a) = {x_a}^k \text{ y } x_k(b) = {x_b}^{k}\]
  con $k\in\{1,2\}$ pero como la deriva es lineal entonces tenemos 
  \begin{equation}\label{eq:6}
    \frac{\partial}{\partial x_k}\left( L + \lambda M\right) - \frac{d}{\,dt}\frac{\partial}{\partial\dot{x_k}}\left( L + \lambda M\right) = 0
  \end{equation}
  asi si consideramos $F = \lambda_0 L + \lambda_1 M$ (o bien $F = L + \lambda M$) entonces la condicion esta dada por 
  \begin{equation}\label{eq:7}
    \frac{\partial}{\partial x_k}F - \frac{d}{\,dt}\frac{\partial}{\partial\dot{x_k}}F = 0 
  \end{equation}
  sujeto a 
  \[x_k(a) = {x_a}^k \text{ y } x_k(b) = {x_b}^{k}\]
  con $k\in\{1,2\}$ 
\end{solution}

\begin{exercise}
  Aplicar el resultado del problema anterior para hallar el valor máximo del funcional
  \[
    J[x_1,x_2] = \int_{a}^{b}(x_1(t)\dot{x_2}(t) - \dot{x_1}(t)x_2(t))\,dt ,
  \]
  donde

  \[
  \begin{aligned}
    & \mathcal{D}_{x_1}\times\mathcal{D}_{x_2} = \{(x_1,x_2)\in\mathcal{C}^1([a,b],\mathbb{R}\times\mathcal{C}^1([a,b],\mathbb{R}): \\
    & \text{ donde }x_k(a) = 0 = x_k(b)\text{ para }k = 1,2\text{ y }G[x_1,x_2] = 1\},
  \end{aligned}
  \]
  con
  \[G[x_1,x_2] = \int_{a}^{b}\sqrt{\dot{x_1}^2(t) + \dot{x_2}^2(t)}\,dt.\]
  ¿Puede establecer alguna relación entre la formulación y la solución de este problema con respecto a los resultados del Ejemplo 5.7?
\end{exercise}
\begin{solution}
  Definamos 
  \[
    \hat{F}[x_1,x_2] = \lambda_0J + \lambda_1\hat{G} = \int_{a}^{b}\left(\lambda_0\left(x_1(t)\dot{x_2}(t) - \dot{x_1}(t)x_2(t)\right) + \lambda_1\left(\sqrt{\dot{x_1}^2(t) + \dot{x_2}^2(t)} - \frac{1}{b-a}\right)\right)\,dt
  \]
  de esto identifiquemos 
  \[
    L = \lambda_0\left(x_1(t)\dot{x_2}(t) - \dot{x_1}(t)x_2(t)\right) + \lambda_1\left(\sqrt{\dot{x_1}^2(t) + \dot{x_2}^2(t)} - \frac{1}{b-a}\right)
  \]
  usando la condición del ejercicio~\ref{ej:4} y suponiendo que $\lambda = \frac{\lambda_1}{\lambda_0}$ se sigue que 
  \[
    \frac{\partial}{\partial x_1}L - \frac{d}{\,dt}\frac{\partial}{\partial\dot{x_1}}F = \dot{x_2} - \frac{d}{\,dt} \left(- x_2 + \lambda\frac{\dot{x_1}}{\sqrt{\dot{x_1}^2 \dot{x_2}^2}}\right) = 2\dot{x_2} - \lambda\frac{d}{\,dt}\left(\frac{\dot{x_1}}{\sqrt{\dot{x_2}^2 + \dot{x_2}^2}}\right) = 0
  \]
  y 
  \[
    \frac{\partial}{\partial x_2}L - \frac{d}{\,dt}\frac{\partial}{\partial\dot{x_2}}F = -\dot{x_1} - \frac{d}{\,dt}\left(x_1 + \lambda \frac{\dot{x_2}}{\sqrt{\dot{x_1}^2 + \dot{x_2}^2}}\right) = -2\dot{x_1} - \lambda\frac{d}{\,dt}\left(\frac{\dot{x_2}}{\sqrt{\dot{x_1}^2 + \dot{x_2}^2}}\right) = 0
  \]
  es decir 
  \begin{equation}\label{eq:8}
    \lambda\frac{d}{\,dt}\left(\frac{\dot{x_1}}{\sqrt{\dot{x_2}^2 + \dot{x_2}^2}}\right) = 2\dot{x_2}
  \end{equation}
\begin{equation}\label{eq:9}
  \lambda\frac{d}{\,dt}\left(\frac{\dot{x_2}}{\sqrt{\dot{x_1}^2 + \dot{x_2}^2}}\right) = -2\dot{x_1}
\end{equation}

ahora consideremos $s$ el parametro de la longitud de arco entonces 
\[\frac{ds}{\,dt} = \sqrt{\dot{x_1}^2 + \dot{x_2}^2} = v\]
luego 
\[
  \begin{aligned}
    & \tau_1 = \frac{dx_1}{\,ds} = \frac{\frac{dx_1}{\,dt}}{\frac{\,dt}{\,ds}} = \frac{\dot{x_1}}{v} \\ 
    & \tau_2 = \frac{dx_2}{\,ds} = \frac{\frac{dx_2}{\,dt}}{\frac{\,dt}{\,ds}} = \frac{\dot{x_2}}{v}
  \end{aligned}
\]
notemos que $\dot{x_2} = v\tau_2$ y $\dot{x_1} = v\tau_1$ sustituyendo en la ecuación~\ref{eq:8} y en la ecuación~\ref{eq:9} se tiene que 
\[\begin{aligned}
  & 2v\tau_2 = \lambda\frac{d}{\,dt}\tau_1 = \lambda\frac{\,ds}{\,dt}\frac{d\tau_1}{\,ds} = \lambda v\frac{d\tau_1}{\,ds} \text{, entonces } \lambda\frac{d\tau_1}{\,ds} = 2\tau_2 \\ 
  & -2v\tau_1 = \lambda\frac{d\tau_2}{\,dt} = \lambda\frac{\,ds}{\,dt}\frac{d\tau_2}{\,ds} = \lambda v\frac{\tau_2}{\,ds}\text{, entonces }\lambda\frac{d\tau_2}{\,ds} = -2\tau_1
\end{aligned}\]
\noindent
simplificando notación $\lambda\dot{\tau_1} = 2\tau_2$ y $\lambda\dot{\tau_2} = -2\tau_1$ ahora si derivamos con respecto a $s$ una vez la primera ecuación tenemos que $\lambda\ddot{\tau_1} = 2\dot{\tau_2}$ y de la segunda se sigue que $\dot{\tau_2} = -\frac{2}{\lambda}\tau_1$.
Entonces sustituyendo se obtiene 
\[\ddot{\tau_1} + \left(\frac{2}{\lambda}\right)^2\tau_1 = 0\]
cuya solución esta dada por 
\[\tau_1 = A\cos(\omega s) + B\sin(\omega s) + C\]
de forma analoga sucede lo mismo con $\tau_2$ asi 
\[\tau_2 = D\cos(\omega s) + F\sin(\omega s) + G\]

por ultimo 
\[J[x_1,x_2]\]
\end{solution}
\end{document}
